\chapter{Sonuç}
\label{ch:sonuc}

EdgeTAM'ın checkpoint'i, Jetson AGX Orin üzerinde dört ayrı TensorRT
motoruna (image encoder, memory attention, memory encoder, SAM head),
fp16 hassasiyetinde ve CUDA Graph tekrar oynatmasıyla
taşınmıştır. Başlangıç hedefleri ölçümle karşılanmıştır: inference
88,30~ms/kareden 26,40~ms/kareye inmiş (\textbf{3,34×}), kuyruk
gecikmesi (p99) 40~ms bütçesinin altında kalmış (39,90~ms), ve fp16
motorları 500 karelik bir klip boyunca fp32 referansına karşı
\textbf{0,9993} ortalama IoU ile eşleşip hata biriktirmemiştir.

Çalışma sırasında ortaya çıkan ve başlangıçta beklenmeyen üç bulgu, bu
raporun asıl katkısını oluşturmaktadır:

\begin{enumerate}
  \item \textbf{Darboğaz image encoder değildi.} Kare başına işlem
        yükünün \%61,9'u memory attention'daydı, bu da kapsamı tek
        modülden EdgeTAM'ın Perceiver ile sıkıştırılmış memory
        mimarisinin tamamına genişletmiştir (\cref{sec:proje-eslesme}).
  \item \textbf{Dört motora ayrışma, gelişigüzel değil ölçülerek
        alınmış bir karardır.} Tek motora birleştirmenin
        kazandırabileceği tek şey ölçülmüş (0,217~ms, karenin
        \%0,5'i) ve eşiğin altında kaldığı için birleştirmemeye karar
        verilmiştir (\cref{sec:fuzyon-karari}).
  \item \textbf{Ön-işleme, ölçülmediğinde modelin kendisi kadar
        pahalıydı.} 1024 çözünürlükte \textasciitilde30~ms'ye kadar
        çıkan bu maliyet iki ayrı sorunun (araştırma amaçlı referans
        kodun float64'e yükseltmesi, ve dosya çözme maliyetinin
        gerçek zamanlı bütçeyle karıştırılması) ayrıştırılmasıyla
        4--5× ucuzlatılmıştır (\cref{sec:onisleme-analizi}).
\end{enumerate}
